\section{Native specializations and prime transport for the normalized cubic tuple}
\label{sec:cubic-native-specialization}

Retain the actual rings and groups of AC-G94--AC-G102, including
the complete coefficient DVR \(\mathcal O\), the closed invariant
ring \(B\subset A\subset R\), the actual tuple
\((\theta_0,\theta_1,\theta_2)=(\chi^j\gamma,1,\chi^j\gamma^{-1})\),
and the surjection
\[
                  \Phi:\mathscr B\longrightarrow B
\tag{G104.1}
\]
from the native ordinary determinant ring.  Write \(H\) for its
native CM Galois group, \(G\) for the actual quotient group, and
\(\pi:H\twoheadrightarrow G\) for AC-G97's map.  Its scalar twist
is \(\kappa\), so the native representation attached to \(r\) is
\(\widetilde r(g)=\kappa(g)r(\pi g)\).
The native full cyclotomic character is denoted by \(\epsilon\),
and \(H'=\ker\epsilon\).  This notation keeps the native convention
of AC-G97; the separation arguments use its kernel and do not
identify it with the duality character \(c_p\).  Corresponding local
groups and their place labels are identified by the actual local
isomorphisms of AC-G97; the \(\zeta_{v,i}\) in (G104.4) use
these identifications.

Let \(C=C_x\) be an actual local completion of AC-G101, with
\(\iota:A\to C\).  At a prime \(P\subset C\), write
\(\operatorname{TI}(P)\) when neither triangular ideal is contained
in \(P\) and at least one odd-weight entry survives, and
\(\operatorname{TE}(P)\) when neither triangular ideal is contained
in \(P\) and all four odd-weight entries vanish.  Thus TE is the
endpoint \(2\)-block plus middle line locus.  Define
\[
\begin{split}
 \widehat{\operatorname{Nice}}(P)\quad\Longleftrightarrow\quad &
 P\text{ prime},\quad \dim C/P=1,\quad p\in P,\\
 &T_-C\not\subset P,\quad T_+C\not\subset P,\quad
 \widehat{\mathfrak c}_{v,ab}\not\subset P
                    \text{ for every }v,\ a\ne b .
\end{split}
\tag{G104.2}
\]
Here the last ideals are the \emph{residually normalized} ideals
of AC-G101.  The distinction from the original source's raw
\texttt{Nice} is essential: a constant nontrivial residual cubic
ratio has no raw equality collision but has finite image.

\begin{theorem}[AC-G104-point: actual native curve points]
\label{AC-G104-point}
For every prime \(P\subset C\) with \(\dim C/P=1\) and \(p\in P\),
there is a finite extension \(k''/k\) and a continuous local
\(\mathcal O\)-algebra homomorphism
\[
 j:C\longrightarrow V=k''[[t]],\qquad \ker j=P,
\tag{G104.3}
\]
for which \(V\) is finite over \(C/P\).  The data
\[
 \rho(g)=\kappa(g)\,j\iota(r(\pi g)),\qquad
 \psi_{v,i}(h)=\kappa(\phi_vh)\,j\iota(\zeta_{v,i}(h))
\tag{G104.4}
\]
form a literal native \texttt{Point}, with symmetric polarization
matrix \(J\).  The composite \(\mathscr B\to B\to C\to V\)
is its native determinant/local-character point map.
The ordered local flag is already defined integrally over \(A\),
before passage to the curve.
\end{theorem}

\begin{proof}
At a coefficient-prime place choose AC-G94's separating element
\(s_v\), and put \(u_i=\zeta_{v,i}(s_v)\in B^\times\).
Their differences are units of \(B\), since their residues are
pairwise distinct.  The labeled projectors are
\[
 E_i=\prod_{h\ne i}\frac{r(s_v)-u_h}{u_i-u_h}\in M_3(A).
\tag{G104.5}
\]
In \(R\), these are the actual rank-one projectors used in the
ordinary defining equations.  Their entries have weights
\(w_a-w_b\), and their diagonal entries belong to \(B\).
Moreover \((E_i)_{ii}\) has residue one and is a unit in \(B\).
Form the matrix \(P_v\) with column
\(E_ie_i/(E_i)_{ii}\).  Every entry belongs to \(A\).
Each monomial in \(\det P_v\) has weight zero, so
\(\det P_v\in B\); its residue is one.  Hence
\(P_v\in\operatorname{GL}_3(A)\).
The lower entries of \(P_v^{-1}r(h)P_v\) vanish and its
diagonal entries are \(\zeta_{v,i}(h)\) in \(R\), by the actual
ordinary equations of AC-G94.  The same identities hold in
\(A\), which is its specified subring, and survive every map
out of \(A\).  This proves the asserted integral flag.
It avoids any assumption that a map out of \(A\) extends to \(R\).

The quotient \(D=C/P\) is a complete Noetherian local domain
of dimension one and characteristic \(p\); primeness gives
\(\varpi\in P\) from \(p\in P\).  Its residue field is finite.
The complete curve normalization theorem gives a finite local
injection into a complete DVR with finite residue field \(k''\).
That DVR is \(k''[[t]]\): its coefficient field is obtained by
the unique simple-root lifts of \(X^{\#k''}-X\), and expansion
in a uniformizer supplies the isomorphism.  This is exactly
\texttt{SelfTwistEigenvalues\_CurveNorm.lean}, theorem
\texttt{curveNorm}; its complete coefficient-field and DVR
argument was read.  Composing with the quotient map proves (G104.3).
The coefficient map from the finite field \(k\) is constant:
its values solve the same simple-root finite-field equations.
Thus it has the native required coefficient tower.

The continuity of the actual representation follows from AC-G95
and locality of the maps into \(C\) and \(V\).  The scalar
\(\kappa\) and \(\pi\) are continuous by AC-G97, so (G104.4)
is continuous.  Polarization and
\(\kappa(g)\kappa(cgc)=\epsilon(g)^{-2}\) give the native
form equation with symmetric invertible matrix \(J\).
The residual representation at \(x\) can be nondiagonal, but
AC-G101 makes it triangular with diagonal \(\boldsymbol\theta\).
Consequently its residual determinant is precisely the required
diagonal residual determinant; no diagonal-matrix reduction is
asserted.  The local character residues are the prescribed ones.
The integral flag just proved gives the required flag over
\(\operatorname{Frac}V\), after the common scalar twist.

For every native away-prime joint residual kernel \(I_u\),
AC-G97 proves \(\pi(I_u)=1\) and \(\kappa|_{I_u}=1\).
Therefore \(\rho(I_u)=1\) exactly.  This checks the native
\texttt{ine} field without a new tame-power or unipotence argument.
All fields of \texttt{Point} have now been checked.
Finally the actual Pair identity of AC-G97, preserved under
base change, identifies the determinant and every local character
of (G104.4) with the indicated composite from \(\mathscr B\).
This is exactly its \texttt{IsPt} assertion.
\end{proof}

\begin{theorem}[AC-G104-nice: normalized nice points are native nice]
\label{AC-G104-nice}
If \(\widehat{\operatorname{Nice}}(P)\), the native point in
AC-G104-point satisfies \texttt{N1} and \texttt{N3}.
At a TI prime its representation is absolutely irreducible
on \(H'\).  At a TE prime it has the native polarized
\(2+1\) decomposition, with absolutely irreducible endpoint
block on \(H'\) and no nonzero equivariant map from that block
to the middle line.  These assertions remain irreducible on
every open subgroup of \(H'\).
If the separately retained native \texttt{N4} input holds,
these are respectively native \texttt{NiceI} and \texttt{NiceE}.
No \texttt{N4} assertion follows from the present geometric proof.
\end{theorem}

\begin{proof}
\emph{1. Simultaneous local separation on open subgroups.}
Put \(N_v=I_v\cap\ker\epsilon\).
By AC-G101 and injectivity of \(C/P\to V\), every pairwise
local-character ratio has infinite image on \(N_v\).
Intersect with the common residual-character kernel to obtain
a finite-index subgroup \(N_v^0\).  On it all ratios are in
\(1+tV\).  Restricting an infinite-image homomorphism to a
finite-index subgroup leaves infinite image, so none becomes
trivial there or on any further open subgroup.
The group \(1+tV\) is torsion free: if \(u^{p^am}=1\)
with \(p\nmid m\), the domain property gives \(u^m=1\), and
\((u-1)(1+u+\cdots+u^{m-1})=0\) has a unit second factor.

For any ambient open subgroup of \(H\), its local preimage
meets \(N_v^0\) in finite index.  The kernels of the finitely
many ratios on this intersection all have infinite index.
They cannot cover the group: the finite coset-cover theorem
would make one of them finite index.  This is the exact
\texttt{CommonFree.lean}, \texttt{commonFree} argument.
Choose an element outside all kernels.  Torsion freeness then
shows that all positive powers of its three eigenvalues are
pairwise distinct.  The scalar \(\kappa\) cancels in ratios.
The element maps into \(H'\).  In particular local characters
are distinct and cannot differ by the native cyclotomic
character, which is one on this element.  This is \texttt{N1}.

\emph{2. Absolute irreducibility before restriction.}
The normalization matrix is
\(Q=r(g_0)-r(g_0^\star)=\operatorname{diag}(a,0,-a)\),
with \(a\in B^\times\), by AC-G102.
Over any field image its three coordinate projectors are
polynomials in \(Q\).  Every invariant subspace is therefore
a sum of coordinate lines.  Draw an arrow \(b\to a\) when
some \(r_{ab}(g)\) is nonzero.  Polarization pairs the two
upper adjacent arrows, and likewise the two lower adjacent
arrows.  An upper adjacent arrow consequently supplies the
chain \(2\to1\to0\), and a lower adjacent arrow supplies
\(0\to1\to2\).  At a TI prime an adjacent chain is present.
Off both triangular loci supplies an arrow in the opposite
direction, either the opposite adjacent chain or the endpoint
arrow completing a cycle.  The directed graph is strongly
connected.  Thus it has no proper nonempty invariant coordinate
subset, over the algebraic closure as well as the original field.

At a TE prime the endpoint arrows in both directions survive,
and the same projector argument makes the endpoint block
absolutely irreducible.  The middle line is a direct summand.
Surjectivity of \(\pi\) and multiplication by a nonzero common
scalar preserve these invariant-subspace assertions on \(H\).

\emph{3. Restriction to the full cyclotomic kernel.}
We apply the precise abstract theorem \texttt{irrNormal} from
\texttt{IrrNormal.lean}, whose full proof, including its
\texttt{CliffPrime} alternative, has been read.  Its inputs are:
an absolutely irreducible representation of prime dimension
over an algebraically closed field; a coefficient domain with
a residual character factorization of every characteristic
polynomial; a local subgroup on which the characteristic
polynomials factor into characters; pairwise local distinction
inside the normal subgroup under consideration; and one
nonscalar residual local witness.  Its conclusion is
irreducibility on that normal subgroup.  The theorem is
abstract in the residual characters; they need not be powers
of one character.

Here the coefficient domain is \(V\), embedded into an
algebraic closure \(\Omega\) of its fraction field.  Reduction
of the characteristic polynomial has the three global residual
characters \(\kappa\theta_i\pi\), since the residual matrix
is triangular.  The ordered local flag supplies the required
local factorization.  For an ambient open normal subgroup
\(W\subset H\), take the normal subgroup \(W\cap H'\).
Step 1 gives its pairwise local distinction.  A tame pro-three
element distinguishes an adjacent residual character from the
middle character, since its residual values before scalar twist
are \((\chi^j,1,\chi^j)\) with \(\chi^j\ne1\).
Every hypothesis of \texttt{irrNormal} in dimension three is
therefore satisfied at a TI prime.

For the endpoint block use the same theorem in prime dimension
two.  The middle local character equals one of the three labeled
characters as a character, not just pointwise: the local ordered
product identity evaluated on the middle line says that if none
of the three character equalities held, choosing one witness
against each would give a product of three nonzero scalars equal
to zero.  Cancel its linear factor from the full characteristic
polynomial.  The remaining two characters give the local block
factorization and are distinguished by Step 1.  The endpoint
residual characters are \(\kappa\chi^j\gamma\) and
\(\kappa\chi^j\gamma^{-1}\).  They are distinct on a local
element with \(\gamma^2\ne1\), one of the actual separating
hypotheses verified in AC-G94.  This witnesses the last input
of \texttt{irrNormal}; the cubic scalar cancels.

Every open subgroup of \(H'\) contains \(W\cap H'\) for some
ambient open normal \(W\): first choose an ambient neighborhood
of one inside its relative open set, then a profinite open normal
subgroup inside that neighborhood.  Irreducibility on the smaller
group proves it on the given open subgroup.  Taking \(W=H\)
also proves the asserted irreducibility on \(H'\) itself.
This completes the changed-residual checks behind the native
\texttt{NiceIrr} and \texttt{NiceIrrE} arguments.

The TE polarization restricts to \(J_2\) and \(J_1\), each
with the same native multiplier \(\epsilon^{-2}\); the multiplier
uses the ambient rank three, exactly as native \texttt{TypeE}
requires.  A map from the irreducible two-dimensional block to
the line has invariant kernel.  A nonzero map would have
one-dimensional kernel, a contradiction.  Thus the native
cross-Hom field is zero.  This verifies every \texttt{TypeE} field.

\emph{4. Regular spectrum in every open coset.}
The full statements and proofs of \texttt{LinClosedOfAdic},
\texttt{StrongIrrOfOpen}, \texttt{RegCoset}, and
\texttt{RegCosetE} apply to these representations.
For clarity their relevant hypotheses are checked here.
A linear equation over \(\Omega\) in matrix entries has closed
zero set on the profinite group: pull its kernel back to the
finite free \(V\)-module of matrices, use separatedness of
the finite module quotient, and the adic continuity of all
entries.  This is \texttt{LinClosedOfAdic}.  Stabilizers of
subspaces are therefore closed.  A finite-index stabilizer is
open, and Step 3 excludes every proper nonzero such invariant
subspace.  Hence the representation, or its endpoint block,
is strongly irreducible in the literal \texttt{StrongIrr} sense.

Let an open set \(U\subset H\) meet \(H'\), and choose
\(s\in U\cap H'\).  Choose an ambient open subgroup \(W\)
with \(sW\subset U\).  Step 1 supplies an element of
\(W\cap H'\) with all eigenvalue ratios of infinite order.
For TI, \texttt{regCoset} in dimension three applies to the
strongly irreducible representation of \(W\cap H'\) and the
matrix \(\rho(s)\).  It gives \(h\in W\cap H'\) for which
\(\rho(sh)\) has separable characteristic polynomial.
For TE, \texttt{regCosetE} applies to the strongly irreducible
endpoint block, its middle character, and the same three
pairwise free eigenvalue ratios.  It gives the identical
separability conclusion for the full \(2+1\) matrix.
Separable characteristic polynomial descends from \(\Omega\)
to \(\operatorname{Frac}V\).  This is precisely native
\texttt{N3}.  Combining it with Steps 1--3 and the separately
assumed \texttt{N4} gives the last assertion.
\end{proof}

\paragraph{The exact thick condition.}
We use the source's algebraic condition with the actual tuple.
A prime \(E\subset C\) is thick if it is TE and, over \(Q=C/E\),
there is a continuous representation \(s\) with the same even
entries as \(r\), with \(s_{10}=s_{12}=0\), satisfying the
ordered local product identity and its reversed starred identity,
and such that after every injective map \(Q\to L\) to a field
there is no vector \(y\) with
\[
 y_1=1,\qquad s(g)y=s_{11}(g)y\quad\text{for every }g.
\tag{G104.6}
\]
This is exactly \texttt{LeAdicClosure.lean},
\texttt{IsThickMod}/\texttt{Thick}; its nonsplitting condition
is retained, not inferred merely from a dimension inequality.
Set
\[
 f:\mathscr B\longrightarrow C,
 \qquad b(P)=f^{-1}(P),\qquad M=\Phi(J)B.
\tag{G104.7}
\]

\begin{theorem}[AC-G104-prime: stable, thick and normalized prime transport]
\label{AC-G104-prime}
A TI prime \(P\subset C\) pulls back to a native stable prime
\(b(P)\).  A thick prime pulls back to a native thick prime.
Assume the retained native assertion that every Point satisfies
\texttt{N4}.  Then the following transport statements hold.

If \(J\) satisfies native \texttt{P1core},
\(\widehat{\operatorname{Nice}}(p)\), \(p\) is TI,
and \(MC\subset p\), then
\[
                 MC\subset P\qquad\text{for every prime }P\subset p.
\tag{G104.8}
\]
If \(J\) satisfies native \texttt{P3core}, let
\(p_0,\ldots,p_h\) be normalized nice TE primes such that
each consecutive pair contains a common thick prime \(E_i\).
If a TI prime \(P_0\subset p_0\) contains \(MC\), then
every TI prime \(P_h\subset p_h\) contains \(MC\).
These are the versions of \texttt{P1'}/\texttt{P23'} with
\texttt{Nice} and \texttt{ThickLink} replaced throughout by
their normalized versions.  The theorem does not assert the
original raw-Nice quantifiers or the final global propagation theorem.
\end{theorem}

\begin{proof}
For TI, apply Step 2 of AC-G104-nice over an algebraic closure
of \(\operatorname{Frac}(C/P)\); that step does not use the
curve or separation hypotheses.  The actual Pair identity gives
an absolutely irreducible realization of the native determinant
at \(b(P)\) in this algebraically closed extension.
The exact generic \texttt{CommutativeAlgebraHeight/StableField.lean},
\texttt{stable\_field}, transports such a realization between
algebraically closed field maps with the same kernel.  Apply it
to the native chosen algebraic closure of the residue field of
\(b(P)\).  This is native stability, with no residual-power
hypothesis added.

Suppose \(E\) is thick, and let \(s\) be its witness.
Twist \(s\pi\) by \(\kappa\).  For every element of the group
algebra, the resulting matrix has zero entries in positions
\((1,0),(1,2)\).  Its endpoint block and middle entry agree
with the block-diagonal representation \(\widetilde r\).
Their characteristic polynomials agree, also after polynomial
scalar extension, and hence their determinant laws agree.
The representation annihilates the native group-algebra
kernel ideal.  Indeed it kills the Cayley--Hamilton generators
by its determinant identity; the local ordered generators by
the two stipulated ordinary identities; and the away inertia
generators and their stars because \(\pi(I_u)=1\) and
\(\kappa(I_u)=1\).  Continuity and the locality of
\(\mathscr B\to Q\) kill the closure as well: reduce modulo
each maximal-ideal power and an open representation kernel,
then intersect all the powers in the separated Noetherian ring.
These are precisely the applied \texttt{th\_CH},
\texttt{th\_closure} and \texttt{kills} arguments in
\texttt{ThickAlg.lean}/\texttt{ThickOf.lean}.

Choose an algebraically closed field \(L\) containing
\(\operatorname{Frac}Q\), and extend the resulting residue-field
embedding of \(b(E)\) to its algebraic closure.  The zero middle
off-row entries make the twisted representation reducible.
It cannot share its determinant law with an absolutely irreducible
three-dimensional representation over the algebraically closed
native residue field.  This is \texttt{ThickAlg.not\_irr}:
Burnside gives all matrix units in the latter representation,
whereas the former's middle-entry character would induce a
unital character of the matrix algebra, which the matrix-unit
relations prohibit.  Thus \(b(E)\) is not stable.

If its native Cayley--Hamilton/ordinary quotient algebra were
semisimple, the invariant endpoint plane in the twisted witness
would have a complementary submodule.  Its projection to the
middle line is an isomorphism, giving a vector \(y_1=1\)
transforming by the middle character.  The exact field-extension
argument is \texttt{ThickAlg.split\_of\_ss}; the middle character
values belong to the algebraically closed native residue field
because they are roots of its determinant characteristic polynomials.
Cancel the nonzero scalar \(\kappa\), and use the surjectivity
of \(\pi\).  This gives (G104.6), contradicting nonsplitting.
The native quotient is therefore not semisimple, proving thickness.

For (G104.8), AC-G104-point produces the native curve with
kernel \(p\).  AC-G104-nice and the retained \texttt{N4}
make it native NiceI.  The ideal containment is equivalently
\[
                 MC\subset Q\quad\Longleftrightarrow\quad J\subset b(Q)
\tag{G104.9}
\]
for any ideal \(Q\subset C\), by the definition of ideal image
and contraction.  The native \texttt{P1core} applied to
\(b(P)\subset b(p)\) gives \(J\subset b(P)\), and (G104.9)
proves the conclusion.

For the chain, construct the native NiceE curve at every
\(p_i\).  Contraction preserves each inclusion
\(E_i\subset p_i,p_{i+1}\).  The first part makes every
\(b(E_i)\) native thick, and \(b(P_0)\), \(b(P_h)\) native
stable.  Formula (G104.9) gives \(J\subset b(P_0)\).
Native \texttt{P3core} applies to this exact finite chain and
gives \(J\subset b(P_h)\); (G104.9) concludes.
This also includes a chain of length zero.  No existence of a
thick chain, seed prime, or prescribed supported curve was used
or asserted.
\end{proof}

\paragraph{Exact sources, open steps and geometric interpretation.}
Every native file cited here is pinned at
\texttt{7cbf7a0c32fca1badf05931879a37c3f8a2b03b4} under
\path{Development/AutomorphicArithmetic/AutomorphicVarieties/}.
The \texttt{Point}, \texttt{N1}--\texttt{N4}, stable/thick and
prime-core definitions are in
\path{WeilGroupCharacters/UniversalOrdinaryDeterminantRing.lean}.
The full containing proofs of \texttt{PointEx}, \texttt{NiceSep},
\texttt{IrrNormal}, \texttt{CliffPrime}, \texttt{CommonFree},
\texttt{NiceIrr}, \texttt{NiceIrrE}, \texttt{NiceTypeE},
\path{NiceSpecializationTwistSatisfiesFirstCondition},
\path{NiceSpecializationTwistSatisfiesThirdCondition},
\texttt{LinClosedOfAdic}, \texttt{StrongIrrOfOpen},
\texttt{RegCoset}, \texttt{RegCosetE}, \texttt{StableOf},
\texttt{ThickAlg}, \texttt{ThickOf},
\texttt{FirstPrimeConditionTransfer} and
\texttt{SecondThirdPrimeConditionTransfer} were read with their
actual residual hypotheses.  These files are in
\path{CurveDVRSpecialization/}; \texttt{stable\_field} has the
separate path given in its application.  Source receipts record
the exact bytes and reading scopes.  No Lean build is claimed.

The proof instantiates the abstract group, linear-algebra and
prime-contraction arguments with the changed tuple and proves
its integral ordinary flag directly.  It does not instantiate
the old \texttt{Setup}, \texttt{Frame.Ok} or \texttt{Frame.Big}
with an incompatible residual character.  In particular the
remaining component argument must use the normalized bad loci
throughout, with the cuts proved in AC-G101.  Parity presentation,
component chaining, the seed comparison and the rank-two corner
needed for global AllGood propagation remain separate until
proved for this tuple.  The arithmetic finite upper certificate,
the original counting budget, prescribed varying-character
support, H04.4 and cubic RH are not proved by this section.

Geometrically, the actual formal representation chart sends
each characteristic-\(p\) curve germ to a native determinant
curve by finite normalization.  The integral projectors supply
its ordered local-system flags.  Removing the normalized finite
ratio subschemes gives the subgroup irreducibility and regular
spectrum needed for native prime transport.  Stable and thick
primes are respectively irreducible determinant fibers and
fibers with a genuinely nonsplit extension module.  This uses
the established finite-etale/local-system and complete-local
dictionary, and makes no new Grothendieck scan attribution.
